%% kapa: exam=1 round=34 year=2023 period=1 form=A subject=economics subject_ko=경제학원론 seq=30 number=70
%% source: Q-Net 감정평가사 기출문제 (cst00309) artlSeq=5211789 file=raw/1008/5211789/2023년도 제34회 감정평가사 1차시험 1교시.pdf
%% conversion: kordoc 4.24.51 -> tools/gen1.py (auto); review=proofread-sonnet; flags=none
\begin{question}{70}
다음과 같은 특징을 가진 화폐시장의 균형에 관한 설명으로 옳지 않은 것은? (단, 폐쇄경제를 가정한다.)
\begin{framedtext}
○ 실질화폐수요는 이자율의 감소함수이다. \par
○ 실질화폐수요는 국민소득의 증가함수이다. \par
○ 명목화폐공급은 중앙은행에 의해 외생적으로 결정된다. \par
○ 물가수준은 단기적으로 고정되어 있으며, 장기적으로 신축적이다. \par
○ 화폐공급이 증가하면 장기적으로 물가수준은 상승한다. \par
\end{framedtext}
\begin{choices}
  \item LM곡선은 우상향한다.
  \item 명목화폐공급이 증가하면 단기적으로 LM곡선이 우측으로 이동한다.
  \item 국민소득이 일정할 때, 명목화폐공급이 이자율에 미치는 영향은 단기보다 장기에서 더 작다.
  \item 실질화폐공급이 증가하면 LM곡선은 우측으로 이동한다.
  \item 장기적으로 실질화폐공급이 변화하지 않는다면, LM곡선은 수직이다.
\end{choices}
\end{question}
\answer{5}
